% Part: counterfactuals
% Chapter: introduction
% Section: material-conditional

\documentclass[../../../include/open-logic-section]{subfiles}

\begin{document}

\olfileid{cnt}{int}{mat}

\olsection{مادي شرطي}

په انګريزۍ کښې د شرطي تر ټولو ساده بڼه د «که \dots نو
\dots» په جوړښت جمله ده؛ د \dots{} پر ځاے خپله جملې راځي،
لکه «که د کور خدمتګار دا کار کړے وي، نو باغوان بې‌ګناه دے»۔
د منطق په لومړنيو کورسونو کښې د $\lif$ نښلوونکي په کارولو د
شرطيو نښه‌ليکنه زده کوو: د \dots له مخې ښودل شوې برخې، مثلاً
په !!{formula}s $!A$ او~$!B$، په نښو وليکئ؛ ټول شرطي بيا
په $!A \lif !B$ ښودل کېږي۔ د ژباړې د نسخې څرګندونه
(OLCNT-001): سرچينه د زده کولو د کلمې لومړے توری پرېږدي؛ د
کورس په سياق کښې د نښه‌ليکنې زده کول مراد دي۔

نښلوونکے $\lif$ \emph{د صدق تابع} دے، يعنې د $!A \lif !B$
د رښتياوالي قيمت---$\True$ يا $\False$---د $!A$ او~$!B$ د
رښتياوالي په قيمتونو ټاکل کېږي: $!A \lif !B$ رښتيا دے هله او
يوازې هله چې $!A$~کاذب يا $!B$~رښتيا وي؛ که داسې نۀ وي،
کاذب دے۔ د رښتياوالي د قيمتونو د ګومارنې~$\pAssign v$ په
نسبت، $\pSat{v}{!A \lif !B}$ تعريفوو هله او يوازې هله چې
$\pSat/{v}{!A}$ يا $\pSat{v}{!B}$ وي۔ د دې معنايي پوهنې
سره نښلوونکے $\lif$ \emph{مادي شرطي} بلل کېږي۔

له دې تعريف څخه څو لومړني منطقي حقيقتونه راوځي۔ تر هر څه
وړاندې، د مادي شرطي لپاره د استنتاج قضيه سمه ده:
\begin{align}
  \text{که } \Gamma, !A \Entails !B \text{ وي، نو } \Gamma \Entails !A \lif !B
\end{align}
هغه د صدق تابع دے: $!A \lif !B$ او $\lnot !A \lor !B$ هم‌ارزښته دي:
\begin{align}
  !A \lif !B & \Entails \lnot !A \lor !B\\
  \lnot !A \lor !B & \Entails !A \lif !B
  \intertext{د مادي شرطي پايله او د هغه د مقدم نفي دواړه دغه
    شرطي په معنايي ډول لازموي:}
  !B & \Entails !A \lif !B\\
  \lnot !A & \Entails !A \lif !B
  \intertext{د مادي شرطي نفي د هغه د مقدم او د پايلې د نفي له
    عطف سره هم‌ارزښته ده: که $!A \lif !B$ کاذب وي،
    $!A \land \lnot !B$ رښتيا دے، او سرچپه هم همداسې ده۔ د
    ژباړې د نسخې څرګندونه (OLCNT-002): سرچينه «کاذب شرطي»
    ليکي؛ په لاندې رسمي استلزامونو کښې د شرطي نفي مراد ده:}
  \lnot(!A \lif !B) & \Entails !A \land \lnot !B\\
  !A \land \lnot !B & \Entails \lnot(!A \lif !B)
  \intertext{مادي شرطي مودوس پوننس مني:}
  !A, !A \lif !B & \Entails !B
  \intertext{د مادي شرطيو پايلې په عطف کښې يوځاے کېدے شي:}
  !A \lif !B,  !A \lif !C & \Entails !A \lif (!B \land !C)
  \intertext{مقدم تل پياوړے کولے شو، يعنې شرطي \emph{يکنواخت} دے:}
  !A \lif !B & \Entails (!A \land !C) \lif !B
  \intertext{مادي شرطي متعدي دے، يعنې د ځنځير قاعده معتبره ده:}
  !A \lif !B, !B \lif !C & \Entails !A \lif !C
  \intertext{مادي شرطي له خپل نقيض عکس سره هم‌ارزښته دے:}
  !A \lif !B & \Entails \lnot !B \lif \lnot !A\\
  \lnot !B \lif \lnot !A & \Entails !A \lif !B
\end{align}

په رياضيکي استدلال کښې دا ټول ګټور او بې‌ستونزې استنتاجونه
دي۔ خو په فلسفي او ژبپوهنې اثارو کښې په غيررياضيکي سياقونو
کښې د ورته استنتاجونو ډېرې ادعا شوې ضدبېلګې شته۔ دا وړانديز
کوي چې د انګريزۍ د «که \dots نو \dots» ويناوو په
نښه‌ليکنه کښې مادي شرطي~$\lif$ مناسب نښلوونکے نۀ دے---يا لږ
تر لږه تل مناسب نۀ دے۔

\end{document}
